% Compile with: uplatex file.tex && dvipdfmx file.dvi
%           or: latexmk -uplatex file.tex
\documentclass[dvipdfmx]{article}
\usepackage{amsmath, amssymb}
\usepackage{tikz}

\title{scenario\_greens\_function\_eom\_commutator}
\date{\today}

\begin{document}
\maketitle

Declare the electron annihilation operator and the scalar coefficients. $\varepsilon$ is the mode energy, $V$ the two-body matrix element.

define c as an electron annihilation operator

\[
  c
\]

define $\varepsilon$ as a scalar parameter

\[
  \epsilon
\]

define V as a scalar parameter

\[
  V
\]

The Hamiltonian. The one-body term is a sum over modes, the two-body term a sum over four indices with an indexed scalar coefficient (which lowers to QuantumAlgebra's `param`). These correspond to $\hat H_0$ (diagonal) and $\hat V$ in the sister scenario.

$H = \sum_n \varepsilon_n c'(n) c(n) + \sum_{p q s n} V_{p,q,s,n} c'(p) c'(q) c(s) c(n)$

\begin{equation}
  H = \sum_n \varepsilon_n c^{\dagger}(n) c(n) + \sum_{p q s n} V_{p,q,s,n} c^{\dagger}(p) c^{\dagger}(q) c(s) c(n)
  \label{eq:E1}
\end{equation}

\textbf{Compute the commutator.} By hand, the one-body term contracts through $\delta_{nm}$ and leaves $-\varepsilon_m c_m$ (the sum disappears); the two-body term gives two $c^\dagger c c$ terms.

commutator of E1 and c(m)

\begin{equation}
  - \varepsilon_{m} {c}_{m} + \sum_{p,q,s} V_{p,m,q,s} {c}_{p}^{\dagger} {c}_{q} {c}_{s} - \sum_{p,q,s} V_{m,p,q,s} {c}_{p}^{\dagger} {c}_{q} {c}_{s}
  \label{eq:E2}
\end{equation}

This is the right-hand side of the Green's function's equation of motion — the step that is written across into the sister scenario. The sign is for $[\hat H,\hat c]$; $[\hat c,\hat H]$ is its negative.

\textbf{From here on the time-carrying operator itself is the object} (PAIR-REP, docs/34 §2). $\hat c(t)=e^{i\hat Ht}\hat ce^{-i\hat Ht}$ is not a new kind of quantity but the sum of nested commutators $\sum_n (it)^n/n!\,\mathrm{ad}_H^n(\hat c)$. With a diagonal $\hat H_0$, $[\hat H_0,\hat c_m]=-\varepsilon_m\hat c_m$ closes on ONE operator, so the sum resums into an exponential \textbf{exactly} — this is not a truncation.

Take the one-body Hamiltonian alone (the interaction does not close the ad-chain; that is slice 2's territory).

$H_0 = \sum_n \varepsilon_n c'(n) c(n)$

\begin{equation}
  H_0 = \sum_n \varepsilon_n c^{\dagger}(n) c(n)
  \label{eq:E3}
\end{equation}

Evolve the annihilation operator. The coefficient $e^{-i\varepsilon_m t}$ is a SCALAR and cannot live in QuantumAlgebra (which decomposes a c-number into cos/sin), so the substance is the (coefficient, operator) PAIR and only the display looks like $\hat c_m(t)$.

evolve c(m) under E3

\begin{equation}
  e^{- i \varepsilon_{m} t} c_{m}
  \label{eq:E4}
\end{equation}

The creation operator too, at a different time $u$. The sign flips ($[\hat H_0,\hat c^\dagger_m]=+\varepsilon_m\hat c^\dagger_m$).

evolve c'(m) in time u under E3

\begin{equation}
  e^{i \varepsilon_{m} u} c^{\dagger}_{m}
  \label{eq:E5}
\end{equation}

\textbf{The different-time anticommutator.} This is what the representation is for: $\{\hat c_m(t),\hat c^\dagger_m(u)\}=f(t)\overline{f(u)}\{\hat c_m,\hat c^\dagger_m\}$ — coefficient product times equal-time algebra. Folding time into an index would have collapsed it to $\delta(t,u)$. The equal-time $\{\hat c,\hat c^\dagger\}=1$ cancels and what remains is a SCALAR with no operator in it: the time dependence of the free propagator.

the anticommutator of E4 and E5

\begin{equation}
  1
  \label{eq:E6}
\end{equation}

\begin{equation}
  e^{i \varepsilon_{m} \left(- t + u\right)}
  \label{eq:E7}
\end{equation}

The commutator keeps its operator ($1-2\hat n_m$), so its result is again a coefficient-times-operator pair.

the commutator of E4 and E5

\begin{equation}
  1 - 2 {c}_{m}^{\dagger} {c}_{m}
  \label{eq:E8}
\end{equation}

\begin{equation}
  e^{i \varepsilon_{m} \left(- t + u\right)} \left(1 - 2 {c}_{m}^{\dagger} {c}_{m}\right)
  \label{eq:E9}
\end{equation}

The anticommutator's result carries no operator, so \textbf{this is the session boundary}: once it is a scalar it can cross to Maxima (SCALAR-HANDOFF). The remaining algebra into $\omega$ space is that side's work.

\textbf{A chain that does not close on ONE operator is still exact if it closes on finitely many.} Try the two-site model with hopping: $\mathrm{ad}$ closes on the two-dimensional span of $\hat c_1$ and $\hat c_2$, and the coefficients become a matrix exponential.

$H_w = w (c'(1) c(2) + c'(2) c(1))$

\begin{equation}
  H_w = w (c^{\dagger}(1) c(2) + c^{\dagger}(2) c(1))
  \label{eq:E10}
\end{equation}

Evolving the annihilation operator now MIXES two operators ($\cos$ and $-i\sin$). That is what differs from the diagonal case, which closes on one.

evolve c(1) under E10

\begin{equation}
  \cos{\left(t w \right)} c_{1} -  i \sin{\left(t w \right)} c_{2}
  \label{eq:E11}
\end{equation}

evolve c'(1) in time u under E10

\begin{equation}
  \cos{\left(u w \right)} c^{\dagger}_{1} + i \sin{\left(u w \right)} c^{\dagger}_{2}
  \label{eq:E12}
\end{equation}

The different-time anticommutator. Two of the four combinations survive, $\cos\cos+\sin\sin$ is collected on the scalar side, and what comes out is a function of the TIME DIFFERENCE alone: $\cos(w(t-u))$ — the free propagator of the two-site model.

the anticommutator of E11 and E12

\begin{equation}
  \cos{\left(w \left(t - u\right) \right)}
  \label{eq:E13}
\end{equation}

\end{document}